% 割权函数次模性的逐边验证：固定一条权为 w 的边 {i,j}，A⊆B，i∉B。
% 按 j 的归属分三种情形比较 Δ(i|A) 与 Δ(i|B)：j∈A：-w,-w；j∈B\A：+w,-w；j∉B：+w,+w。
\begin{tikzpicture}[line width=.8pt,font=\figurefont\fontsize{8}{10}\selectfont,
  vtx/.style={circle,draw=FigureStroke,fill=FigureBlueFill,minimum size=6mm,inner sep=0pt,line width=.8pt}]
  \foreach \k/\jx/\jy/\ga/\gb/\ttl in {0/0.55/0.75/{-w}/{-w}/{$j\in A$},1/1.75/0.75/{+w}/{-w}/{$j\in B\setminus A$},2/3.25/0.75/{+w}/{+w}/{$j\notin B$}}{
    \begin{scope}[xshift=\k*36mm]
      \draw[FigureBlue,fill=FigureBlueFill,rounded corners=3pt] (0,0) rectangle (1.2,1.5);
      \draw[FigureBlue,rounded corners=3pt,line width=1pt] (-0.15,-0.15) rectangle (2.35,1.65);
      \node[FigureBlue,anchor=south west,inner sep=1pt] at (0.02,1.5) {$A$};
      \node[FigureBlue,anchor=south west,inner sep=1pt] at (2.0,1.65) {$B$};
      \node[vtx,fill=FigureYellowFill] (i) at (2.75,-0.6) {$i$};
      \node[vtx] (j) at (\jx,\jy) {$j$};
      \draw[FigureStroke,line width=1.1pt] (i)--(j);
      \node[FigureInk,anchor=north] at (1.25,-1.15) {\ttl};
      \node[FigureInk,anchor=north,align=center] at (1.25,-1.6) {$\Delta(i\mid A)=\ga$\\$\Delta(i\mid B)=\gb$};
    \end{scope}
  }
\end{tikzpicture}
